\chapter{Program Context and Objective}

\section{Operational Problem}

Assembly is a time-dependent process. The operator may hold a part out of view,
repeat a motion, pause, or perform a valid operation in a different order from
another operator. A useful vision system therefore has to interpret a sequence
of observations rather than classify isolated images. It must also make a
decision from the video available so far; using future frames would make an
offline result look stronger than a real deployment can be.

The project objective was to build that complete chain: capture representative
egocentric video, label the procedure, evaluate candidate recognition approaches,
adapt the strongest feature representation to causal inference, and expose the
result through interfaces suitable for demonstration.

\section{Development Objectives}

\begin{enumerate}
    \item Establish a repeatable procedure-step recognition workflow using the
    IndustReal research setting.
    \item Collect and label an in-house egocentric dataset on the turbofan
    engine model.
    \item Compare model architectures, visual representations, temporal
    modeling, and decoding choices using consistent metrics.
    \item Train and validate a streaming model that does not inspect future
    frames.
    \item Demonstrate the result through replay, live video, and a wearable
    HoloLens interface.
\end{enumerate}

\section{Scope Boundaries and Success Criteria}

The project was scoped as a procedure-state assistant, not as a complete
factory inspection system. The target behavior was to identify the part or
procedure state visible in the egocentric stream, preserve the order and
duration of the recognized steps, and expose a stable result with enough
latency and robustness for a live demonstration. It was not intended to
replace an approved work instruction, determine torque or fit, certify a
completed installation, or independently authorize an operator to continue.

Progress was evaluated at three levels. Research experiments had to show a
measurable improvement in temporal recognition quality. The own-dataset phase
had to produce repeatable labels, held-out evaluation, and stress cases that
challenged memorized procedure order. The final demonstration had to connect
the model to a functioning replay path, a live stream, and a wearable
presentation. This staged definition prevented a high offline score from being
treated as sufficient evidence of operational readiness.

\section{Development Path}

\begin{table}[H]
\centering
\small
\begin{tabularx}{\textwidth}{@{}p{0.22\textwidth}X@{}}
\toprule
\textbf{Stage} & \textbf{Outcome} \\
\midrule
PSR research & Baselines and controlled experiments on IndustReal established
the useful combination of motion and appearance information. \\
Own data collection & Forty curated egocentric recordings were captured on a
ten-part engine model, covering assembly and disassembly. \\
Annotation and stress testing & Dense completion labels were produced and
additional recordings were generated for scrambled, reversed, partial, and
wrong-step behavior. \\
Streaming conversion & Offline sequence modeling was converted to a causal
temporal model with a fixed-lag decoder. \\
Operational demonstration & Replay and live services were connected to a
browser interface and a HoloLens overlay. \\
\bottomrule
\end{tabularx}
\caption{Project progression from research to demonstration.}
\label{tab:development-path}
\end{table}

\section{Transition from Research to the Own Dataset}

The IndustReal experiments selected the feature and temporal-modeling direction;
they did not replace the need for an internal dataset. The engine-model
recordings were collected to test the selected approach under the camera
perspective, part vocabulary, lighting, operator behavior, and procedure
variability relevant to the intended demonstration. This transition is
important: the project moved from comparing architectures on a benchmark to
building a complete data-to-deployment loop in which recording, annotation,
training, evaluation, replay, and live presentation could all be controlled by
the project team.

\section{Assembly Object}

\begin{figure}[H]
\centering
\includegraphics[width=0.72\textwidth]{engine_model_reference.png}
\caption{The turbofan engine model used for the project dataset and
demonstration.}
\label{fig:engine-model}
\end{figure}

The model provides ten named parts and supports both installation and removal
procedures. The visual appearance of a part is only one cue: the system must
also use motion, temporal context, and the state of the procedure to determine
which step is taking place.

\section{System Boundary at the Demonstration Handoff}

\begin{table}[H]
\centering
\small
\begin{tabularx}{\textwidth}{@{}X X@{}}
\toprule
\textbf{Within the demonstrated scope} & \textbf{Outside the demonstrated
scope} \\
\midrule
Recognizing a procedure state from egocentric video. &
Certifying torque, fit, or final installation quality. \\
Committing stable steps and showing procedure progress. &
Replacing approved work instructions or human sign-off. \\
Flagging out-of-order or attention-requiring activity. &
Providing a complete fault taxonomy for every assembly failure. \\
Delivering state to replay, live, and wearable interfaces. &
Spatially anchoring every instruction to a physical part. \\
\bottomrule
\end{tabularx}
\caption{Capability boundary used to interpret the current demonstration.}
\label{tab:scope-boundary}
\end{table}
